\documentclass[11pt]{article}
\usepackage[T1]{fontenc}
\usepackage[utf8]{inputenc}
\usepackage{lmodern,amsmath,amssymb,amsthm,mathtools,booktabs,array}
\usepackage[margin=27mm]{geometry}
\usepackage{microtype,hyperref,xcolor}
\definecolor{inkplum}{HTML}{54384D}
\hypersetup{colorlinks=true,linkcolor=inkplum,citecolor=inkplum,urlcolor=inkplum,
 pdftitle={Fixed-size subset sums: sharp support, Renyi entropy and tail bounds},
 pdfauthor={Anonymous}}
\newtheorem{theorem}{Theorem}[section]
\newtheorem{lemma}[theorem]{Lemma}
\newtheorem{proposition}[theorem]{Proposition}
\newtheorem{corollary}[theorem]{Corollary}
\theoremstyle{remark}
\newtheorem{remark}[theorem]{Remark}
\newcommand{\R}{\mathbb R}
\newcommand{\Z}{\mathbb Z}
\newcommand{\Prb}{\mathbb P}
\newcommand{\E}{\mathbb E}
\newcommand{\supp}{\operatorname{supp}}
\newcommand{\Hyp}{\operatorname{Hyp}}
\newcommand{\Ber}{\operatorname{Ber}}
\newcommand{\rank}{\operatorname{rank}}
\newcommand{\one}{\mathbf 1}
\setlength{\emergencystretch}{2em}
\allowdisplaybreaks[1]
\title{Fixed-size subset sums:\\ sharp support, R\'enyi entropy and tail bounds}
\author{Anonymous}
\date{14 November 2026 {\small(version 0.1.0-candidate, DOI \href{https://doi.org/10.5281/zenodo.24241302}{10.5281/zenodo.24241302})}}
\begin{document}
\maketitle
\vspace{-1.5em}
\begin{abstract}
Choose \(k\) of \(n\) given real numbers uniformly at random and add them.
This paper determines exactly how concentrated such a fixed-size subset sum
can be, in three senses. First, for a nonzero vector with zero total, at least
\(2\binom{n-2}{k-1}\) of its \(k\)-subsets have a nonzero sum whenever
\(k\ge5\) and \(n\ge2k+1\), with equality only for \((1,-1,0,\ldots,0)\) up to
scaling and reordering. This proves a 2018 conjecture of Mogilnykh,
Vorob'ev and Valyuzhenich on the smallest support of an eigenfunction in the
first nontrivial eigenspace of the Johnson graph, which they had checked by
computer for \(n\le600\). A formula with all equality cases is given for every
\(n\) and \(k\), and a missing small equality case is added. Second, the
minimum R\'enyi entropy of the sum is found for every finite order, answering
a question of Ai, Ku\v{z}elka and Pelekis: a Bernoulli law is extremal except
at half-sampling, where a three-point law takes over above an explicit order,
so the universal Bernoulli bound holds exactly for orders up to two. Third,
for the two-valued populations to which Aydinian and Blinovsky conjecturally
reduce the Manickam--Mikl\'os--Singhi problem, the probability of a
nonnegative sum is minimised by a single outlier whenever \(n\ge4k\); this
extends the known range \(n\ge8k\) to the threshold of that conjecture.
Analogues over Abelian groups and a small-ball inequality complete the
picture.
\end{abstract}

\tableofcontents

\section{Introduction}\label{sec:intro}

\subsection{The question}
Take \(n\) real numbers \(x_1,\ldots,x_n\), choose \(k\) of them uniformly at
random without replacement, and add them. Throughout, \(n\ge2\) and
\(1\le k\le n-1\). Write
\[
 \Omega_{n,k}=\{S\subseteq[n]:|S|=k\},\qquad
 L_x(S)=\sum_{i\in S}x_i,\qquad X=L_x(S)\ \text{with }S\text{ uniform on }\Omega_{n,k}.
\]
How concentrated can the random sum \(X\) be? Unlike a sum of independent
terms, a fixed-size sample has rigid dependence: the sample of the
complementary indices has sum \(\sum_ix_i-X\). The question has three
classical forms, which this paper takes up in turn.

\begin{enumerate}
\item \emph{The zero atom.} If \(\sum_ix_i=0\) and \(x\ne0\), how many
\(k\)-subsets can have sum exactly zero? This is the minimum-support problem
for the first nontrivial eigenspace of the Johnson graph (Section~\ref{subsec:known}).
\item \emph{The whole law.} How small can the Shannon or R\'enyi entropy of
\(X\) be when \(x\) is not constant? This is the ``entropic'' version of the
Manickam--Mikl\'os--Singhi problem introduced by Ai, Ku\v{z}elka and Pelekis
\cite{AKP}.
\item \emph{The sign.} If \(\sum_ix_i=0\), how small can \(\Prb(X\ge0)\) be?
This is the Manickam--Mikl\'os--Singhi (MMS) problem itself.
\end{enumerate}

A single principle links the three. Every population is dominated, in a
precise majorization sense, by a population that takes only two values; the
sum of a two-valued population is an affine image of a hypergeometric count.
For the zero atom and for the whole law (up to majorization) this reduction is
proved below, and each question becomes one about a few atoms of
hypergeometric laws; the paper's new ingredients are sharp estimates for those
atoms. For the sign the reduction to two-valued populations is a conjecture of
Aydinian and Blinovsky, and the paper solves the resulting hypergeometric
problem.

\subsection{What was known}\label{subsec:known}

\paragraph{Minimum supports in the Johnson graph.}
The Johnson graph \(J(n,k)\) has the \(k\)-subsets of \([n]\) as vertices,
two being adjacent when they share \(k-1\) elements. For an eigenvalue of a
graph one may ask for the smallest number of vertices on which a nonzero
eigenfunction can be supported; this ``uncertainty principle'' question has
been studied for many families of graphs, and is surveyed by Sotnikova and
Valyuzhenich \cite{SV}. For \(J(n,k)\), the eigenspace of the largest
nontrivial eigenvalue \(k(n-k)-n\) consists exactly of the functions \(L_x\)
with \(\sum_ix_i=0\) \cite{VMV,MVV}; indeed
\(\sum_{T\sim S}L_x(T)=\{k(n-k)-n\}L_x(S)\), and \(x\mapsto L_x\) is
injective on zero-sum vectors (if all \(k\)-sums agree, all coordinates
agree, and the zero total makes them vanish). So the first question asks for
\[
 m(n,k)=\min\bigl\{\,|\{S\in\Omega_{n,k}:L_x(S)\ne0\}|\;:\;x\in\R^n\setminus\{0\},\ \textstyle\sum_ix_i=0\bigr\}.
\]
Vorob'ev, Mogilnykh and Valyuzhenich \cite{VMV} found the minimum support of
every fixed Johnson eigenspace when \(n\) is larger than an unspecified
\(n_0(k)\). Mogilnykh, Vorob'ev and Valyuzhenich \cite{MVV} showed that a
minimiser is either a multiple of the coordinate difference
\(D=(1,-1,0,\ldots,0)\) or a two-valued vector, so that \(m(n,k)\) is the
smaller of two explicit expressions. They compared the two expressions by
computer for \(6\le n\le600\), recorded the balanced value and three small
exceptions, and stated \cite[Conjecture~1]{MVV}:
\[
 m(n,k)=2\binom{n-2}{k-1}\qquad(k\ge5,\ n\ge2k+1).
\]
What remained open was therefore a uniform proof of a binomial comparison,
valid for every \(n\), together with the absence of further exceptions for
\(k\le4\) beyond the computed range.

\paragraph{Entropy of subset sums.}
Ai, Ku\v{z}elka and Pelekis \cite{AKP} proved the sharp Shannon bound
\(H_1(X)\ge H_1(\Ber(k/n))\), where \(\Ber(k/n)\) is the Bernoulli law with
parameter \(k/n\), for nonconstant zero-sum populations. Their argument (their
Lemma~3) compares the law of \(X\) by majorization with hypergeometric laws
\(\Hyp(n,i,k)\), the law of the number of marked items in a random
\(k\)-sample from \(n\) items of which \(i\) are marked (Section~\ref{subsec:notation}).
For R\'enyi entropy their Corollary~5 reduces the problem to minimising
\(H_\alpha(\Hyp(n,i,k))\) over \(1\le i<n\). They wrote that
they were unable to solve this optimisation and that their numerical
experiments suggest the minimum is at \(i=1\) \cite[Section~5]{AKP}. An earlier paper of the same authors \cite{AKP25} uses
majorization of populations to obtain tail bounds for sampling without
replacement.

\paragraph{The sign problem.}
The MMS conjecture asserts that if \(n\ge4k\) and \(\sum_ix_i=0\), then
\(\Prb(X\ge0)\ge k/n\), i.e.\ at least \(\binom{n-1}{k-1}\) of the
\(k\)-subsets have nonnegative sum \cite{MS88}; the population
\((n-1,-1,\ldots,-1)\) shows that \(k/n\) would be sharp. It is known for
\(n\ge10^{46}k\) \cite{Pok}. Aydinian and Blinovsky \cite{AB} conjectured that
for every \(n\) and \(k\) the two-valued populations are extremal, which
reduces the problem to minimising the hypergeometric tails
\(\Prb(\Hyp(n,i,k)\ge ik/n)\) over \(1\le i\le n-1\). Ai and Pelekis
\cite{AP} solved this tail problem for \(n\ge8k\) and remarked that they were
unable to reach \(n\ge4k\).

\paragraph{Other value groups.}
The affine version of the zero-atom question, with values in an arbitrary
Abelian group, is a minimum-distance question for codes on the Boolean slice.
Amireddy, Behera, Srinivasan and Sudan \cite{ABSS} proved a degree-one
distance bound \((r-1)/n\), \(r=\min(k,n-k)\), as part of a polynomial
distance lemma for low-degree testing; the binary case is the
Dumer--Kapralova theory of spherically punctured biorthogonal codes
\cite{DK,BCIM}.

\subsection{Notation}\label{subsec:notation}
Binomial coefficients with an invalid lower index are zero, and
\(r=\min(k,n-k)\).
\begin{itemize}
\item \(\Hyp(n,i,k)\) is the number of marked items in a uniform \(k\)-sample,
without replacement, from \(n\) items of which \(i\) are marked; it has mass
\(\binom ij\binom{n-i}{k-j}/\binom nk\) at \(j\) and mean \(ik/n\). If the
population takes the value \(y_+\) on \(i\) indices and \(y_-<y_+\) on the
rest, then \(X=ky_-+(y_+-y_-)\Hyp(n,i,k)\).
\item \(\Ber(p)\) denotes the two-point probability vector \((1-p,p)\). For
\(k\ge2\) put
\[
 U=\Bigl(\tfrac12,\tfrac12\Bigr),\qquad
 q_k=\frac{k}{2k-1},\qquad
 Q_k=\Bigl(q_k,\ \tfrac{1-q_k}2,\ \tfrac{1-q_k}2\Bigr)
 =\Bigl(\tfrac{k}{2k-1},\tfrac{k-1}{2(2k-1)},\tfrac{k-1}{2(2k-1)}\Bigr).
\]
When \(n=2k\), \(U\) is the law of the sample sum of a population with one
outlier, \((1,0,\ldots,0)\), and \(Q_k\) is the law for two outliers,
\((1,1,0,\ldots,0)\).
\item Probability vectors are compared after sorting in decreasing order and
padding with zeros: \(p\preceq q\) (\(q\) \emph{majorizes} \(p\)) if the sum
of the \(\ell\) largest entries of \(p\) is at most that of \(q\) for every
\(\ell\). R\'enyi entropy is taken in bits:
\[
 H_\alpha(p)=\frac{\log_2\sum_jp_j^\alpha}{1-\alpha}\ (0<\alpha<\infty,\
 \alpha\ne1),\qquad H_1(p)=-\sum_jp_j\log_2p_j .
\]
Every \(H_\alpha\), \(0<\alpha<\infty\), is strictly Schur-concave on padded
vectors: \(p\preceq q\) implies \(H_\alpha(p)\ge H_\alpha(q)\), with equality
only when \(p\) and \(q\) agree after sorting. (This follows from strict
convexity or concavity of \(u\mapsto u^\alpha\), and from strict concavity of
\(-u\log u\) at \(\alpha=1\).) A comparison \(p\preceq q\) is \emph{proper} if
\(p\) and \(q\) do not agree after sorting.
\end{itemize}

\subsection{Main results}

\begin{theorem}[First Johnson eigenspace: complete support formula]\label{thm:support}
For all \(n\ge2\) and \(1\le k\le n-1\),
\[
 m(n,k)=
 \begin{cases}
 2,&r=1,\\
 2\binom{n-2}{r-2},&n=2r,\ r\ge2,\\
 6,&(n,r)=(6,2),\\
 39,&(n,r)=(9,3),\\
 110,&(n,r)=(10,4),\\
 2\binom{n-2}{r-1},&\text{otherwise}.
 \end{cases}
\]
All minimising vectors, up to nonzero scalar multiples and permutations of
coordinates, are listed below (superscripts denote repetition).
\begin{center}
\begin{tabular}{@{}lll@{}}
\toprule
Parameters & Minimising vectors & Status of the value\\
\midrule
\(r=1\) & \(D\) & elementary\\
\(n=2r,\ r\ge2\) & \(((r-1)^2,(-1)^{2r-2})\) & \cite[Section 5]{MVV}\\
\((n,r)=(5,2)\) & \(D\) and \((1,1,-1,-1,0)\) & second vector new\\
\((n,r)=(6,2)\) & \((1^3,(-1)^3)\) & \cite{MVV}\\
\((n,r)=(8,2)\) & \(D\) and \((1^4,(-1)^4)\) & \cite{MVV}\\
\((n,r)=(9,3)\) & \((2^3,(-1)^6)\) & \cite{MVV}\\
\((n,r)=(10,4)\) & \((1^5,(-1)^5)\) & \cite{MVV}\\
all other \(n>2r\ge4\) & \(D\) & computed in \cite{MVV} for \(r\ge3\), \(n\le600\)\\
\bottomrule
\end{tabular}
\end{center}
In particular \(m(n,k)=2\binom{n-2}{k-1}\) for \(k\ge5\) and \(n\ge2k+1\),
which is \cite[Conjecture~1]{MVV}.
\end{theorem}

The extra vector \((1,1,-1,-1,0)\) at \((n,k)=(5,2)\) is three-valued, so it is
absent from the two-class description in \cite[Theorem~1]{MVV} and its
restatement \cite[Theorem~5]{SV}; Remark~\ref{rem:52} explains why. The value
\(m(5,2)=6\) is unaffected.

For the remaining questions, \(x\) is any nonconstant real population; its
mean need not be zero. The next theorem describes the law of \(X\) completely
up to majorization, and is the source of the entropy results.

\begin{theorem}[Majorization frontier]\label{thm:frontier}
Let \(x\in\R^n\) be nonconstant and let \(p\) be the mass vector of \(X\).
\begin{enumerate}
\item If \(n\ne2k\), or \(n=2\), then \(p\preceq\Ber(k/n)\).
\item If \(n=2k\ge4\), then \(p\preceq U\) or \(p\preceq Q_k\).
\end{enumerate}
All of these vectors are attained, and at \(n=2k\ge4\) the vectors \(U\) and
\(Q_k\) are incomparable. Consequently the minimum over nonconstant \(x\) of any
symmetric Schur-concave function of the mass vector, invariant under padding
by zeros, is its value at \(\Ber(k/n)\), respectively the smaller of its
values at \(U\) and \(Q_k\).
\end{theorem}

\begin{theorem}[Sharp R\'enyi entropy]\label{thm:entropy}
For every \(0<\alpha<\infty\),
\[
 \min_{x\ \text{nonconstant}} H_\alpha(X)=
 \begin{cases}
 H_\alpha(\Ber(k/n)),& n\ne2k,\ \text{or }(n,k)=(2,1),\\
 \min\{1,H_\alpha(Q_k)\},& n=2k\ge4.
 \end{cases}
\]
For each \(k\ge2\) there is a unique \(\alpha_k>1\) with
\begin{equation}\label{eq:transition}
 \frac12\Bigl(\frac{2k}{2k-1}\Bigr)^{\alpha_k}
 +\Bigl(\frac{k-1}{2k-1}\Bigr)^{\alpha_k}=1 .
\end{equation}
At \(n=2k\ge4\) the minimum is \(1=H_\alpha(U)\) for \(\alpha<\alpha_k\) and
\(H_\alpha(Q_k)<1\) for \(\alpha>\alpha_k\). One has \(\alpha_2=2\),
\(\alpha_k\) increases strictly, and
\(\alpha_k=(2k-\tfrac12)\ln2-\tfrac{\ln2}{24k}+O(k^{-2})\).
Every minimising population takes exactly two values. Their multiplicities
are \(\{1,n-1\}\) when \(n\ne2k\) and when \(n=2k\), \(\alpha<\alpha_k\);
they are \(\{2,n-2\}\) when \(n=2k\), \(\alpha>\alpha_k\); at
\(\alpha=\alpha_k\) exactly these two families minimise. (For \((n,k)=(2,1)\)
every nonconstant population has multiplicities \(\{1,1\}\) and law \(U\).)
\end{theorem}

\begin{corollary}[The universal Bernoulli bound]\label{cor:universal}
The inequality \(H_\alpha(X)\ge H_\alpha(\Ber(k/n))\) holds for every
nonconstant real population and all \(n,k\) if and only if \(0<\alpha\le2\).
\end{corollary}

So the numerical suggestion in \cite{AKP} is correct away from half-sampling
and, at half-sampling, exactly for \(\alpha\le\alpha_k\). For example, at
\(n=4\), \(k=2\) the population \((1,1,0,0)\) has law \((1/6,2/3,1/6)\),
whose third-order power sum \(11/36\) exceeds \(1/4\); hence
\(H_3=\tfrac12\log_2(36/11)<1=H_3(\Ber(1/2))\).

\begin{theorem}[Hypergeometric tails from \(n=4k\)]\label{thm:tail}
Let \(k\ge1\) and \(n\ge4k\). For every \(1\le i\le n-1\),
\[
 \Prb\bigl(\Hyp(n,i,k)\ge ik/n\bigr)\ \ge\ \frac kn ,
\]
with equality at \(i=1\); when \(n>4k\), equality holds only at \(i=1\).
\end{theorem}

Theorem~\ref{thm:tail} answers the question left open in \cite{AP} for
\(4k\le n<8k\). Combined with the conjecture of Aydinian and Blinovsky it would
give the MMS conjecture in its full range (Corollary~\ref{cor:AB}). Below
\(n=4k\) the single outlier is no longer extremal for some \(n\)
(Remark~\ref{rem:below4k}).

\begin{theorem}[Abelian groups]\label{thm:groups}
Let \(G\) be a nontrivial Abelian group and let
\(f(S)=c+\sum_{i\in S}a_i\) with \(c,a_1,\ldots,a_n\in G\), not identically
zero on \(\Omega_{n,k}\). Then the proportion of \(k\)-subsets with
\(f(S)\ne0\) is at least
\[
 \delta_G(n,k)=
 \begin{cases}
 r/n,&n\ne2k,\\
 (k-1)/(2k-1),&n=2k\ge4,\\
 1/2,&(n,k)=(2,1),
 \end{cases}
\]
and this is attained. Moreover, for nonconstant populations
\(a_1,\ldots,a_n\in G\) and \(n=2k\ge4\), the minimum of \(H_\alpha(X)\)
is \(H_\alpha(\Ber(q_k))\) if \(G\) has an element of order two, and
\(\min\{1,H_\alpha(Q_k)\}\) otherwise; for \(n\ne2k\) it is
\(H_\alpha(\Ber(k/n))\).
\end{theorem}

\subsection{What is new, and what is not}
\begin{itemize}
\item Theorem~\ref{thm:support}: the reduction to two classes of vectors is
\cite[Theorem~1]{MVV}; the balanced value and the small exceptions are in
\cite[Section~5]{MVV}. New are a uniform proof of the comparison (the
integer-mean atom estimate, Lemma~\ref{lem:meanatom}), hence
\cite[Conjecture~1]{MVV}; the absence of further exceptions for \(k\le4\);
and the extra equality case at \((5,2)\).
\item Theorems~\ref{thm:frontier} and \ref{thm:entropy}: the reduction of
arbitrary populations to hypergeometric laws is \cite[Lemma~3]{AKP}; the
balanced maximum atom \(q_k\) is in \cite[Section~5]{MVV}; the link between
majorization and R\'enyi entropy is classical \cite{MWW}. New are the
frontier at half-sampling, the evaluation of the R\'enyi optimisation with the
threshold \(\alpha_k\), all equality cases, and a parity majorization chain
for balanced hypergeometric laws (Theorem~\ref{thm:parity}).
\item Theorem~\ref{thm:tail} is new. Its proof rests on recent results of Ai,
Karagiannis, Ku\v{z}elka and Pelekis \cite{AKKP}: a lemma on random variables
with increasing failure rate (Lemma~\ref{lem:AKKP}) and its consequence
\cite[Theorem~3(1)]{AKKP} for hypergeometric laws with non-integer mean. The
new steps are the reduction to these results for \(n\ge4k\), the
integer-mean case (via the identity used for Theorem~\ref{thm:support}) and the
case of small mean.
\item Theorem~\ref{thm:groups}: the binary distance is due to Dumer and
Kapralova \cite{DK} (see \cite[Theorem~1.2]{BCIM}); the balanced value for
\(n\ge8\) can be extracted from the proof of \cite[Theorem~6.1]{ABSS}. New
are the uniform exact formula, which sharpens the stated bound \((r-1)/n\) of
\cite{ABSS} (by \(1/n\) away from half-sampling, by \((k-1)/(2k(2k-1))\) at
\(n=2k\)) and removes its condition \(n\ge8\), and the entropy
dichotomy governed by elements of order two.
\end{itemize}
The paper also gives a small-ball inequality for sampling without replacement
(Corollary~\ref{cor:smallball}), and recovers from the atom bound the known
largest families of \(k\)-sets whose incidence matrix has rank below \(n\)
(Corollary~\ref{cor:rank}).

\subsection{Method and organisation}
Section~\ref{sec:support} proves the conjectured range of
Theorem~\ref{thm:support}. A count due to \cite{MVV} handles vectors with three
or more distinct values (Lemma~\ref{lem:three}); for two-valued vectors the
zero atom is a hypergeometric atom at an integer mean, and
Lemma~\ref{lem:meanatom} bounds it by comparing it with its four neighbours.
Section~\ref{sec:small} completes Theorem~\ref{thm:support} for \(k\le4\) and
for \(n=2k\). Section~\ref{sec:maj} proves Theorem~\ref{thm:frontier}: a chain
construction reduces every population to two-valued ones
(Lemma~\ref{lem:chain}), an atom bound settles \(n\ne2k\), and a parity
majorization chain settles \(n=2k\). Section~\ref{sec:entropy} derives
Theorem~\ref{thm:entropy} and Corollary~\ref{cor:universal} and classifies
equality. Section~\ref{sec:conc} gives the largest-atom, small-ball and
incidence-rank consequences. Section~\ref{sec:tail} proves
Theorem~\ref{thm:tail}. Section~\ref{sec:abelian} treats Abelian groups and
proves Theorem~\ref{thm:groups}. Section~\ref{sec:open} lists open problems.

\section{The conjectured range}\label{sec:support}

This section proves Theorem~\ref{thm:support} for \(k\ge5\), \(n\ge2k+1\).
Complementation \(S\mapsto[n]\setminus S\) negates zero-sum \(k\)-sums into
\((n-k)\)-sums, so \(m(n,k)=m(n,n-k)\) and we may assume \(k\le n/2\).

\begin{lemma}[Three distinct values; \cite{MVV}]\label{lem:three}
Let \(G\) be an Abelian group, \(c\in G\), and let \(a\in G^n\) have at least
three distinct coordinates. Then
\[
 |\{S\in\Omega_{n,k}:c+L_a(S)\ne0\}|\ \ge\ 2\binom{n-2}{k-1}.
\]
If \(G=\R\), \(c=0\), \(\sum_ia_i=0\), \(k+1\le n-3\), and equality holds, then
\(a\) is a scalar multiple of a permutation of \(D\).
\end{lemma}
\begin{proof}
Let \(a_1,a_2,a_3\) be distinct. For each \((k-1)\)-subset \(A\) of
\(\{4,\ldots,n\}\), the three values \(c+L_a(A)+a_\ell\), \(\ell\le3\), are
distinct, so at least two are nonzero. For each \((k-2)\)-subset \(A\) of
\(\{4,\ldots,n\}\), the three values \(c+L_a(A)+a_\ell+a_{\ell'}\),
\(\ell<\ell'\le3\), are distinct by cancellation in \(G\), so again at least
two are nonzero. These classes of \(k\)-sets are disjoint, and
\(2\binom{n-3}{k-1}+2\binom{n-3}{k-2}=2\binom{n-2}{k-1}\). This counting step
is the reduction in \cite[Theorem~1]{MVV}.

For equality in the real zero-sum case, every \(k\)-set inside
\(\{4,\ldots,n\}\) must have sum zero. Since \(1\le k<n-3\), exchanging one
element of such a set shows that \(a_4,\ldots,a_n\) are all equal, and a
vanishing \(k\)-sum makes them zero. Equality in the class with one special
index then says that exactly two of \(a_1,a_2,a_3\) are nonzero; the total is
zero, so they are opposite.
\end{proof}

\paragraph{Two-valued vectors and cell counts.}
Let \(x\) be zero-sum and two-valued, with the larger value on \(i\) indices,
\(1\le i\le n-1\). Then \(L_x(S)\) is an affine function of
\(J=|S\cap I|\sim\Hyp(n,i,k)\), where \(I\) is the set of those indices, and
\(L_x(S)=0\) exactly when \(J=ik/n\). So a zero sum is possible only when
\(j=ik/n\) is an integer, and then
\begin{equation}\label{eq:meanatom}
 \Prb(L_x(S)=0)=p_j=\frac{\binom ij\binom{n-i}{k-j}}{\binom nk}.
\end{equation}
In that case write the four cell counts of the \(2\times2\) table (marked or
not, sampled or not) at the mean as
\[
 a=j,\qquad b=i-j,\qquad c=k-j,\qquad d=n-i-k+j .
\]
The mean condition \(nj=ik\) is equivalent to \(ad=bc\), and adjacent atoms
satisfy
\begin{equation}\label{eq:ratios}
 \frac{p_{j-1}}{p_j}=\frac{ad}{(b+1)(c+1)}=\frac{bc}{(b+1)(c+1)},\qquad
 \frac{p_{j+1}}{p_j}=\frac{bc}{(a+1)(d+1)}=\frac{ad}{(a+1)(d+1)} .
\end{equation}
The vector \(D\) has \(2\binom{n-2}{k-1}\) nonzero \(k\)-sums, a proportion
\(2k(n-k)/(n(n-1))\). Put
\begin{equation}\label{eq:T}
 T(n,k)=1-\frac{2k(n-k)}{n(n-1)},
\end{equation}
the proportion of zero \(k\)-sums of \(D\). A two-valued vector beats \(D\)
exactly when its zero atom \(p_j\) exceeds \(T(n,k)\).

\begin{lemma}[Integer-mean atom]\label{lem:meanatom}
If \(k\ge5\), \(n\ge2k+1\), \(1\le i<n\) and \(j=ik/n\) is an integer, then
\(p_j<T(n,k)\).
\end{lemma}
\begin{proof}
Since \(p_j\) is unchanged by \(i\mapsto n-i\) (with \(j\mapsto k-j\)), we may
assume \(i\le n/2\). Then \(a\le c\) because \(j=ik/n\le k/2\); \(a<b\)
because \(j=ik/n<i/2\) (as \(k<n/2\)); and \(c<d\) because \(ad=bc\) and
\(a<b\). Also
\(a\ge1\), since \(i\ge1\) and \(j=ik/n\) is a positive integer. Hence all four
counts are positive integers, and \(a+c=k\).

\emph{Case \(a\ge2\).} Then \(b\ge a+1\ge3\); \(c\ge k/2>2\), so \(c\ge3\);
and \(d\ge c+1\ge4\). Since \(a,b,c,d\ge2\), the neighbouring atoms
\(j\pm1\) and \(j\pm2\) all carry positive mass. The functions
\(t\mapsto t/(t+1)\) and \(t\mapsto t^2/((t+1)(t+2))\) increase, so
\eqref{eq:ratios} gives
\[
 \frac{p_{j-1}}{p_j}=\frac{b}{b+1}\cdot\frac{c}{c+1}\ge\frac34\cdot\frac34=\frac9{16},
 \qquad
 \frac{p_{j+1}}{p_j}=\frac{a}{a+1}\cdot\frac{d}{d+1}\ge\frac23\cdot\frac45=\frac8{15}.
\]
For the two-step ratios, using \(ad=bc\) once more,
\begin{align*}
 \frac{p_{j-2}}{p_j}
 &=\frac{b^2}{(b+1)(b+2)}\cdot\frac{c^2}{(c+1)(c+2)}
   \Bigl(1-\frac1a\Bigr)\Bigl(1-\frac1d\Bigr)
 \ge\frac{9}{20}\cdot\frac9{20}\cdot\frac12\cdot\frac34=\frac{243}{3200}>\frac1{16},\\
 \frac{p_{j+2}}{p_j}
 &=\frac{a^2}{(a+1)(a+2)}\cdot\frac{d^2}{(d+1)(d+2)}
   \Bigl(1-\frac1b\Bigr)\Bigl(1-\frac1c\Bigr)
 \ge\frac13\cdot\frac{8}{15}\cdot\frac23\cdot\frac23=\frac{32}{405}>\frac1{16}.
\end{align*}
Since the five atoms have total mass at most one,
\[
 1>p_j\Bigl(1+\frac9{16}+\frac8{15}+\frac1{16}+\frac1{16}\Bigr)=p_j\cdot\frac{533}{240},
 \qquad p_j<\frac{240}{533}<\frac5{11}.
\]
For fixed \(k\), \(T(n,k)\) increases in \(n\ge2k+1\): the derivative of its
rational extension has numerator \(n^2-2kn+k>0\). Hence
\(T(n,k)\ge T(2k+1,k)=k/(2k+1)\ge5/11\).

\emph{Case \(a=1\).} Then \(c=k-1\ge4\), \(b\ge2\), \(d=bc\), and
\(n=a+b+c+d=(b+1)k\ge3k\). The immediate neighbours suffice:
\[
 \frac{p_{j-1}}{p_j}=\frac{b}{b+1}\cdot\frac{c}{c+1}\ge\frac23\cdot\frac45=\frac8{15},\qquad
 \frac{p_{j+1}}{p_j}=\frac12\cdot\frac{d}{d+1}\ge\frac12\cdot\frac89=\frac49,
\]
so \(p_j(1+8/15+4/9)<1\) and
\[
 p_j<\frac{45}{89}<\frac{11}{21}\le\frac{5k-3}{9k-3}=T(3k,k)\le T(n,k).
\]
The strict gaps are \(5/11-240/533=25/5863\) and \(11/21-45/89=34/1869\).
\end{proof}

\begin{proof}[Proof of Theorem~\ref{thm:support} for \(k\ge5\), \(n\ge2k+1\)]
A one-valued zero-sum vector is zero. Vectors with three or more distinct
values have at least \(2\binom{n-2}{k-1}\) nonzero \(k\)-sums by
Lemma~\ref{lem:three}, with equality only for \(D\) since
\(n-3\ge2k-2\ge k+1\). A two-valued vector with \(ik/n\notin\Z\) has no zero
\(k\)-sum at all; with \(ik/n\in\Z\), Lemma~\ref{lem:meanatom} shows it has
strictly more nonzero \(k\)-sums than \(D\).
\end{proof}

\section{The remaining parameters}\label{sec:small}

This section completes the proof of Theorem~\ref{thm:support}, again with
\(k\le n/2\). The argument for \(n>2k\) is the same as in
Section~\ref{sec:support}: three-valued vectors are handled by
Lemma~\ref{lem:three}, and a two-valued competitor with an integer mean has
zero atom \(p_j\), to be compared with \(T(n,k)\). The cell counts
\(a,b,c,d\) are as in Section~\ref{sec:support}, with \(a\le c\).

\paragraph{\(k=4\).}
If \(a=1\), \(c=3\), then \(b\ge2\), \(d=3b\), \(n=4(b+1)\ge12\), and
\[
 \frac{p_{j-1}}{p_j}\ge\frac12,\qquad
 \frac{p_{j+1}}{p_j}\ge\frac37,\qquad
 \frac{p_{j+2}}{p_j}=\frac{b(b-1)}{(3b+1)(3b+2)}\ge\frac1{28};
\]
the last bound follows from writing the ratio as
\(\frac{b-1}{b}\cdot\frac{b^2}{(3b+1)(3b+2)}\), a product of increasing
factors. It gives \(p_j\le28/55<17/33=T(12,4)\le T(n,4)\).
If \(a=c=2\), then \(b=d\ge3\), \(n=2b+4\), and
\[
 p_j=\frac{3(b+2)(b+1)}{2(2b+3)(2b+1)} ,
\]
which decreases strictly in \(b\) (the ratio of consecutive terms is
\((b+3)(2b+1)/((b+1)(2b+5))<1\)). At \(b=3\), \(n=10\), it gives
\(p_j=10/21\) and support \(110<112\): the exception \(m(10,4)=110\). For
\(b\ge4\), \(p_j\le5/11<17/33\le T(n,4)\).

\paragraph{\(k=3\).}
Here \(a=1\), \(c=2\), \(d=2b\), \(n=3(b+1)\), with \(b\ge2\). At \(b=2\),
\(n=9\), the zero atom is \(15/28\), giving support \(39<42\): the exception
\(m(9,3)=39\). For \(b\ge3\) the immediate ratios are at least \(1/2\) and
\(3/7\), so \(p_j\le14/27<13/22=T(12,3)\le T(n,3)\).

\paragraph{\(k=2\).}
An integer mean requires \(a=c=1\), \(b=d\), so \(n=2b+2\) is even and
\(p_j=n/(2(n-1))\). Then
\[
 T(n,2)-p_j=\frac{(n-2)(n-8)}{2n(n-1)},
\]
which gives the exception \(m(6,2)=6\), the tie \(m(8,2)=12\) (both \(D\) and
\((1^4,(-1)^4)\) minimise), and strict superiority of \(D\) for even \(n\ge10\).
For odd \(n\) no two-valued competitor exists. For \(n\ge6\) the hypothesis
\(k+1\le n-3\) of Lemma~\ref{lem:three} holds, which gives the remaining
equality cases (at \(n=6\) the two-valued vector is strictly better, so \(D\)
is not a minimiser there).

At \(n=5\), \(k=2\), Lemma~\ref{lem:three} gives \(m(5,2)\ge6\), and \(D\)
attains \(6\), but its equality statement does not apply since \(n-3=k\). A
vector attains \(6\) exactly when four of its ten pairs have sum zero. If
\(z\) coordinates vanish, the zero pairs are the \(\binom z2\) pairs of zeros
and the pairs of opposite nonzero values. The cases \(z=4,5\) are impossible
for a nonzero zero-sum vector. If \(z=3\), the two nonzero coordinates are
opposite, giving \(D\). If \(z=2\), there is no zero pair among the zeros
beyond one, and the three nonzero coordinates (with zero total) contain no
opposite pair, so there is one zero pair only. If \(z=1\), four opposite pairs
among four nonzero values require two copies each of \(u\) and \(-u\), giving
\((1,1,-1,-1,0)\). If \(z=0\), four or more opposite pairs among five values
\(\pm u\) need multiplicities \((1,4)\), \((2,2)\) or \((2,3)\) for some
\(u\); the \((2,2)\) pattern leaves one further nonzero coordinate and a
nonzero total, and the other two patterns also have nonzero total; values
from two different absolute levels give at most three opposite pairs. Hence
the minimisers at \((5,2)\) are \(D\) and \((1,1,-1,-1,0)\).

\paragraph{\(k=1\).}
A nonzero zero-sum vector has at least two nonzero coordinates, with equality
exactly for multiples of \(D\).

\paragraph{\(n=2k\), \(k\ge2\).}
Three distinct values give at least \(2\binom{2k-2}{k-1}\) nonzero sums by
Lemma~\ref{lem:three}, more than the claimed value. A two-valued vector has a
zero atom only when \(i=2s\) is even, and then the atom is
\[
 c_s=\frac{\binom{2s}{s}\binom{2k-2s}{k-s}}{\binom{2k}{k}},\qquad1\le s\le k-1 .
\]
The sequence is symmetric under \(s\mapsto k-s\), and
\[
 \frac{c_{s+1}}{c_s}=\frac{(2s+1)(k-s)}{(s+1)(2k-2s-1)},
\]
whose denominator minus numerator is \(k-2s-1\). So \((c_s)\) decreases up to
the middle and then increases; its maximum is \(c_1=c_{k-1}=k/(2k-1)=q_k\),
attained only at \(s\in\{1,k-1\}\). The minimum support is therefore
\(\binom{2k}{k}-2\binom{2k-2}{k-1}=2\binom{2k-2}{k-2}\), with the listed
minimisers. This value and construction are recorded in
\cite[Section~5]{MVV}. This completes the proof of Theorem~\ref{thm:support}.

\begin{remark}[The equality case at \((5,2)\)]\label{rem:52}
The vector \((1,1,-1,-1,0)\) has six nonzero pair sums and is optimal for
\(J(5,2)\), but it is neither a coordinate difference nor two-valued. It is
omitted from \cite[Theorem~1]{MVV} because the proof there uses injectivity of
an inclusion map on \(n-3\) coordinates, which fails when \(n-3=k\); the same
classification is repeated in \cite[Theorem~5]{SV}. An exhaustive search over
all candidate vectors for \(n\le7\) (see ``Data, verification and AI use'' at the end) confirms that
this is the only omission in that range.
\end{remark}

\section{Majorization by hypergeometric laws}\label{sec:maj}

This section proves Theorem~\ref{thm:frontier}. Lemma~\ref{lem:chain} reduces
any population to two-valued ones; Lemma~\ref{lem:atom} then settles
\(n\ne2k\), and the parity chain of Theorem~\ref{thm:parity} settles \(n=2k\).

\begin{lemma}[Chain reduction; \cite{AKP}]\label{lem:chain}
Suppose the sorted population has a strict gap after its \(i\)-th smallest
entry, so that \(i\) lower entries lie below \(n-i\) upper entries, and let
\(g>0\) be the difference between the smallest upper and the largest lower
entry. Then there is a coupling of \(S\) with a random chain
\(S_0,S_1,\ldots\) of \(k\)-sets such that \(S=S_J\) with
\(J\sim\Hyp(n,n-i,k)\) independent of the chain, and the sums
\(L_x(S_j)\) increase in \(j\) by at least \(g\) at each step. Consequently
the mass vector of \(X\) is majorized by that of \(\Hyp(n,i,k)\).
\end{lemma}
\begin{proof}
Order the lower group and the upper group independently and uniformly at
random. For each feasible \(j\), \(\max(0,k-i)\le j\le\min(k,n-i)\), let
\(S_j\) consist of the first \(j\) upper and the first \(k-j\) lower entries.
Passing from \(S_j\) to \(S_{j+1}\) replaces a lower entry by an upper one, so
the sum increases by at least \(g\). Choose \(J\sim\Hyp(n,n-i,k)\)
independently. Given \(J=j\), the set \(S_J\) is uniform among \(k\)-sets with
\(j\) upper entries, so \(S_J\) is uniform on \(\Omega_{n,k}\).

Fix any \(\ell\) real values. Each chain takes each of them at most once, so it
hits them at at most \(\ell\) ranks \(j\), and the conditional probability of
hitting them is at most the sum of the \(\ell\) largest masses of
\(\Hyp(n,n-i,k)\). Averaging over chains and maximising over the \(\ell\)
values gives the majorization. Finally \(\Hyp(n,n-i,k)\) and \(\Hyp(n,i,k)\)
have the same masses in reverse order.
\end{proof}
This is the chain mechanism of \cite[Lemma~3 and Section~4]{AKP}, stated at an
arbitrary strict gap; translating the population removes the need for a zero
mean. The increment bound \(g\) is used in Corollary~\ref{cor:smallball}.

\begin{lemma}[Atoms away from half-sampling]\label{lem:atom}
If \(1\le k<n/2\) and \(1\le i<n\), then for every integer \(j\)
\[
 \binom ij\binom{n-i}{k-j}\le\binom{n-1}{k}.
\]
Equality holds exactly at \((i,j)=(1,0)\) and \((n-1,k)\), and, when
\(n=2k+1\), also at \((2,1)\) and \((n-2,k-1)\).
\end{lemma}
\begin{proof}
We may assume the left side is positive. Put \(\ell=k-j\) and \(m=n-i\). Since
\(j+\ell=k<n/2=(i+m)/2\), interchanging the two groups if necessary we may
assume \(\ell<m/2\), so that \(\binom m{\ell+1}\ge\binom m\ell\). Vandermonde's
identity gives
\begin{align*}
 \binom{n-1}{k}
 &=\sum_h\binom{i-1}{h}\binom m{k-h}
 \ge\binom{i-1}{j}\binom m\ell+\binom{i-1}{j-1}\binom m{\ell+1}\\
 &\ge\Bigl(\binom{i-1}{j}+\binom{i-1}{j-1}\Bigr)\binom m\ell
 =\binom ij\binom m\ell .
\end{align*}
For equality in this orientation: if \(j=0\), the omitted terms vanish only
when \(i=1\); if \(m=1\), then \(\ell=0\), which gives the other endpoint.
Assume \(m\ge2\) and \(j\ge1\). If \(j\ge2\), the omitted term \(h=j-2\) is
positive because \(\ell+2\le m\); hence \(j=1\). Equality in
\(\binom m{\ell+1}\ge\binom m\ell\) forces \(m=2\ell+1\). If \(i=1\), then
\(n=1+m=2\ell+2=2k\), excluded by \(k<n/2\); so \(i\ge2\). If \(i\ge3\), the
omitted term \(h=2\) is positive; thus \(i=2\) and \(n=2k+1\). Undoing the
interchange gives the list.
\end{proof}

Only the first partial sum matters when comparing with a two-point vector, so
Lemma~\ref{lem:atom} gives
\begin{equation}\label{eq:unbalancedmaj}
 \Hyp(n,i,k)\preceq\Ber(k/n)\qquad(k<n/2,\ 1\le i<n).
\end{equation}

\begin{theorem}[Parity majorization at half-sampling]\label{thm:parity}
Fix \(k\ge2\) and let \(h_t\) be the mass vector of \(\Hyp(2k,t,k)\),
\(1\le t\le2k-1\). Then
\[
 h_{t+2}\preceq h_t\qquad(1\le t\le k-1).
\]
The comparison is proper for \(t<k-1\), and \(h_{k+1}\) and \(h_{k-1}\) agree
after sorting. Consequently \(h_t\preceq U\) for odd \(t\) and \(h_t\preceq Q_k\)
for even \(t\).
\end{theorem}
\begin{proof}
Centre the support with \(\rho=2j-t\):
\[
 \eta_t(\rho)=\Prb\bigl(\Hyp(2k,t,k)=(t+\rho)/2\bigr)
 =\frac{\binom{t}{(t+\rho)/2}\binom{2k-t}{(2k-t-\rho)/2}}{\binom{2k}{k}},
\]
supported on \(\rho\in\{-t,-t+2,\ldots,t\}\) when \(t\le k\). The law is
symmetric in \(\rho\) (exchange the sampled and unsampled halves) and unimodal,
because hypergeometric masses are log-concave in \(j\); and \(\eta_t=\eta_{2k-t}\).
For \(t\le k-1\) the laws \(\eta_t\) and \(\eta_{t+2}\) live on the same parity
lattice, and on the common support \(|\rho|\le t\),
\begin{equation}\label{eq:radialratio}
 \frac{\eta_{t+2}(\rho)}{\eta_t(\rho)}
 =\frac{(t+2)(t+1)}{(2k-t)(2k-t-1)}\cdot
   \frac{(2k-t)^2-\rho^2}{(t+2)^2-\rho^2}.
\end{equation}
For \(t<k-1\) we have \(t+2<2k-t\), so the right side increases strictly in
\(\rho^2\). At the two new points \(\rho=\pm(t+2)\), \(\eta_t=0<\eta_{t+2}\).
Both laws have total mass one, so the difference
\(\eta_t-\eta_{t+2}\), as a function of \(|\rho|\), changes sign exactly once,
from positive to negative. Hence for every \(R\) the central mass
\(\sum_{|\rho|\le R}\eta_t(\rho)\) is at least the corresponding mass of
\(\eta_{t+2}\): before the sign change this is immediate, and after it the
central difference equals minus a tail difference, which is nonnegative.

By symmetry and unimodality, the \(\ell\) largest masses of each law can be
taken from the centre outwards. If \(t\) is even and \(\ell\) is odd, or \(t\)
is odd and \(\ell\) is even, they form a full central block
\(\{|\rho|\le R\}\), and the inequality above applies; both laws have the same
parity, so the blocks match. For the remaining \(\ell\) the sum of the \(\ell\) largest masses is the
average of the sums for \(\ell-1\) and \(\ell+1\), because the two masses of a
symmetric pair are equal; so these inequalities follow as well. Thus
\(h_{t+2}\preceq h_t\), properly for \(t<k-1\) since \(\eta_{t+2}\) has larger
support. At \(t=k-1\), \(h_{k+1}=h_{k-1}\) after sorting, by \(\eta_t=\eta_{2k-t}\).
Iterating, and using \(\eta_t=\eta_{2k-t}\) for \(t>k\), every odd \(t\)
reduces to \(t=1\), where \(h_1=U\), and every even \(t\) to \(t=2\), where
\(h_2=Q_k\) by the computation of \(c_1\) in Section~\ref{sec:small}.
\end{proof}

\begin{proof}[Proof of Theorem~\ref{thm:frontier}]
A nonconstant population has a strict gap, so Lemma~\ref{lem:chain} gives
\(p\preceq\Hyp(n,i,k)\) for some \(1\le i<n\). If \(k<n/2\),
\eqref{eq:unbalancedmaj} gives \(p\preceq\Ber(k/n)\); if \(k>n/2\), apply this
to the complementary sample, whose sum is \(\sum_ix_i-X\) and has the same
mass vector, and note \(\Ber((n-k)/n)\) and \(\Ber(k/n)\) agree after sorting.
For \(n=2\) the law is \(U=\Ber(1/2)\). If \(n=2k\ge4\),
Theorem~\ref{thm:parity} gives \(p\preceq U\) or \(p\preceq Q_k\). The
populations \((1,0,\ldots,0)\) and, at \(n=2k\), \((1,1,0,\ldots,0)\) attain
\(\Ber(k/n)\), \(U\) and \(Q_k\). Finally \(q_k>1/2\), so \(Q_k\not\preceq U\),
while the two largest entries of \(U\) sum to one, so \(U\not\preceq Q_k\).
The last sentence of the theorem follows from Schur-concavity.
\end{proof}

Related orderings of symmetric hypergeometric laws concern normal
approximation \cite{MSch} or peakedness about a fixed centre as the sample size
varies \cite{Levin}; Theorem~\ref{thm:parity} instead compares balanced laws
centred at their own, moving, means.

\begin{remark}[A shorter route to the frontier value]\label{rem:shortroute}
At half-sampling the frontier can also be reached from the known maximal atom
\(q_k\) \cite[Section~5]{MVV} and the symmetry \(X\overset d=\sum_ix_i-X\).
If the largest atom is at most \(1/2\), the law is majorized by \(U\).
Otherwise the largest atom is the centre of symmetry, with mass
\(c\le q_k\), and the second largest is at most \((1-c)/2\), so the two
largest sum to at most \((1+q_k)/2\); this is majorization by \(Q_k\).
Theorem~\ref{thm:parity} is stronger: it orders all balanced hypergeometric
laws within each parity class, and its strictness is what the equality
classification in Section~\ref{sec:entropy} uses.
\end{remark}

\section{Entropy thresholds and equality cases}\label{sec:entropy}

This section proves Theorem~\ref{thm:entropy} and
Corollary~\ref{cor:universal}. By Theorem~\ref{thm:frontier}, the minimum is
\(H_\alpha(\Ber(k/n))\) when \(n\ne2k\) and \(\min\{H_\alpha(U),H_\alpha(Q_k)\}\)
when \(n=2k\ge4\), with \(H_\alpha(U)=1\). It remains to locate the crossover
and to classify the minimisers.

\paragraph{The threshold \(\alpha_k\).} Put
\[
 u_k=\frac{2k}{2k-1},\qquad v_k=\frac{k-1}{2k-1},\qquad
 F_k(s)=\tfrac12u_k^s+v_k^s .
\]
Since \(Q_k=(u_k/2,v_k/2,v_k/2)\), the sum of the \(s\)-th powers of its entries
is \(2^{-s}(u_k^s+2v_k^s)=2^{1-s}F_k(s)\), so
\[
 H_s(Q_k)=1+\frac{\log_2F_k(s)}{1-s}\qquad(s\ne1).
\]
Thus \(H_s(Q_k)<1\) exactly when \(F_k(s)>1\) for \(s>1\), and exactly when
\(F_k(s)<1\) for \(s<1\). The function \(F_k\) is strictly convex, with
\(F_k(0)=3/2\), \(F_k(1)=1\) and \(F_k(s)\to\infty\). Its derivative at one is
\(F_k'(1)=\ln2-H_1(Q_k)\ln2<0\): indeed \(Q_k\preceq Q_2\), and
\(H_1(Q_2)\ln2=\ln3-\tfrac13\ln2>\ln2\) since \(27>16\). So \(F_k\) has exactly
one root \(\alpha_k>1\) besides \(s=1\); \(F_k>1\) on \((0,1)\), \(F_k<1\) on
\((1,\alpha_k)\) and \(F_k>1\) on \((\alpha_k,\infty)\). Equation
\eqref{eq:transition} is \(F_k(\alpha_k)=1\). At \(s=1\),
\(H_1(Q_k)\ge H_1(Q_2)>1\). This gives the two regimes of
Theorem~\ref{thm:entropy}. Direct substitution gives \(\alpha_2=2\). Both the
largest entry and the sum of the two largest entries of \(Q_k\) decrease
strictly with \(k\), so \(Q_{k+1}\prec Q_k\) properly; strict
Schur-concavity at \(s=\alpha_k\) gives \(H_{\alpha_k}(Q_{k+1})>1\), hence
\(F_{k+1}(\alpha_k)<1\) and \(\alpha_{k+1}>\alpha_k\).

\begin{proof}[Proof of Corollary~\ref{cor:universal}]
If \(n\ne2k\) the bound holds at every order by Theorem~\ref{thm:entropy}. If
\(n=2k\ge4\), the minimum equals \(1=H_\alpha(\Ber(1/2))\) exactly when
\(\alpha\le\alpha_k\); since \(\alpha_k\ge\alpha_2=2\), the bound holds for all
\(n,k\) when \(\alpha\le2\), and for \(\alpha>2\) it fails at \(n=4\), \(k=2\).
\end{proof}

\begin{remark}[Size of \(\alpha_k\)]\label{rem:asymp}
Let \(A_k=\ln2/\ln u_k\), so \(u_k^{A_k}=2\) and \(F_k(A_k)=1+v_k^{A_k}>1\);
hence \(\alpha_k<A_k\). The inequality \(-\ln(1-x)\le x/(1-x)\) with
\(x=1/(2k)\) gives \(A_k\ge(2k-1)\ln2\); as \(v_k<1/2\),
\(v_k^{A_k-1}\le2^{1-(2k-1)\ln2}\), which decays exponentially in \(k\). So,
for all sufficiently large \(k\),
\(F_k(A_k-1)=1/u_k+v_k^{A_k-1}=1-\frac1{2k}+v_k^{A_k-1}<1\), and
\(\alpha_k\in(A_k-1,A_k)\). On that interval
\(F_k'(s)=\tfrac12u_k^s\ln u_k+v_k^s\ln v_k\ge\tfrac1{u_k}\ln u_k-v_k^{A_k-1}|\ln v_k|\),
which is at least \(\tfrac12\ln u_k\) for all sufficiently large \(k\), since
\(\ln u_k\sim1/(2k)\) and \(1/u_k\to1\) while \(v_k^{A_k-1}|\ln v_k|\) decays
exponentially. The
mean value theorem applied to \(F_k(A_k)-F_k(\alpha_k)=v_k^{A_k}\) gives
\(0<A_k-\alpha_k\le2v_k^{A_k}/\ln u_k=O\bigl(k\,e^{-2(\ln2)^2k}\bigr)\), since
\(v_k^{A_k}=2^{-A_k}(2v_k)^{A_k}\le2^{-A_k}\) and \(A_k\ge(2k-1)\ln2\). Expanding
\(A_k=\ln2/(-\ln(1-1/(2k)))=\ln2\,(2k-\tfrac12-\tfrac1{24k}+O(k^{-2}))\) proves
the asymptotic formula in Theorem~\ref{thm:entropy}. Rounded values are
\(\alpha_3\approx3.594270\), \(\alpha_4\approx5.089878\),
\(\alpha_5\approx6.531142\) and \(\alpha_{10}\approx13.512604\). Neither the
expansion nor these values is used elsewhere.
\end{remark}

\paragraph{Equality cases.} We prove the last sentence of
Theorem~\ref{thm:entropy}. Fix \(0<\alpha<\infty\) and a minimising population
\(x\); by strict Schur-concavity, its mass vector \(p\) agrees after sorting
with the frontier vector it is majorized by.

\emph{A two-point law forces two values.} Suppose \(X\) takes exactly two
values, differing by \(d>0\). Comparing two \(k\)-sets that share \(k-1\)
elements shows \(x_\ell-x_{\ell'}\in\{0,\pm d\}\) for all \(\ell,\ell'\); three
distinct population values would have extreme difference \(2d\), so there are
exactly two, with multiplicities \(i\) and \(n-i\). Their sum takes
\(\min(i,n-i,k,n-k)+1\) values, so for \(2\le k\le n-2\) we get
\(i\in\{1,n-1\}\); for \(k\in\{1,n-1\}\), equality of the masses with
\(\Ber(k/n)\) gives the same multiplicities.

\emph{Unbalanced case \(n\ne2k\).} Here \(p=\Ber(k/n)\) is two-point, and the
previous paragraph gives multiplicities \(\{1,n-1\}\).

\emph{Balanced case below the threshold, \(\alpha<\alpha_k\).} Choose any strict
gap of \(x\), with \(t\) entries below it. If \(t\) were even,
Lemma~\ref{lem:chain} and Theorem~\ref{thm:parity} would give
\(p\preceq Q_k\) and \(H_\alpha(p)\ge H_\alpha(Q_k)>1\), contradicting
minimality. So \(t\) is odd and \(p\preceq U\) with equality, so \(p=U\) is
two-point and \(x\) has multiplicities \(\{1,n-1\}\).

\emph{Balanced case above the threshold, \(\alpha>\alpha_k\).} Now
\(H_\alpha(U)=1>H_\alpha(Q_k)\), so no gap may have odd \(t\), and \(p=Q_k\)
after sorting. For an even gap with \(t\notin\{2,2k-2\}\), the chain from
\(t\) down to \(2\) (or from \(2k-t\) down to \(2\)) in Theorem~\ref{thm:parity}
contains a proper step, so \(p\prec Q_k\) properly, which is impossible. Hence
every gap has \(t\in\{2,2k-2\}\): either \(x\) has two values with
multiplicities \(\{2,2k-2\}\), or three values \(y_1<y_2<y_3\) with
multiplicities \(2,2k-4,2\) (for \(k\ge3\)). In the latter case fix a set
\(M\) of \(k-2\) of the middle entries (possible since \(2k-4\ge k-2\)), and
note that \(2k-4-(k-2)=k-2\ge1\) middle entries remain. Adding to \(M\) the
pairs \(\{y_1,y_1\}\), \(\{y_1,y_2\}\), \(\{y_1,y_3\}\), \(\{y_2,y_3\}\),
\(\{y_3,y_3\}\) gives \(k\)-sets whose sums
\[
 2y_1<y_1+y_2<y_1+y_3<y_2+y_3<2y_3
\]
(each plus \(L_x(M)\)) are five distinct values, whereas \(Q_k\) has only three
atoms. So three values are impossible. For \(k=2\) the only even gap is \(t=2\),
which also forces two values.

\emph{At the threshold, \(\alpha=\alpha_k\).} Both \(U\) and \(Q_k\) give
entropy one. If some gap has odd \(t\), then \(p\preceq U\) with equal entropy,
so \(p=U\) and the one-outlier family results. If all gaps are even, the
previous paragraph applies and gives the two-outlier family. These families
are distinct, since \(q_k>1/2\). This completes the proof of
Theorem~\ref{thm:entropy}.

\section{Largest atoms, small balls and incidence ranks}\label{sec:conc}

This section records three consequences of Section~\ref{sec:maj}: the largest
possible atom, its extension to short intervals, and an extremal problem for
incidence matrices.

\begin{corollary}[Largest atom]\label{cor:concentration}
For a nonconstant real population,
\[
 \sup_t\Prb(X=t)\le
 \begin{cases}
 \max(k,n-k)/n,&n\ne2k,\\
 q_k=k/(2k-1),&n=2k\ge4,\\
 1/2,&n=2.
 \end{cases}
\]
Every bound is attained. Away from half-sampling, equality holds exactly for
two-valued populations with multiplicities \(\{1,n-1\}\) and, when
\(n=2r+1\), also \(\{2,n-2\}\). At half-sampling with \(n\ge4\), equality
requires multiplicities \(\{2,n-2\}\).
\end{corollary}
\begin{proof}
The bounds are the first partial sums in Theorem~\ref{thm:frontier}. For
equality away from half-sampling assume \(k<n/2\). A population with three
distinct values has, for every \(t\), at most a proportion \(T(n,k)\) of
\(k\)-sums equal to \(t\), by Lemma~\ref{lem:three}; and
\(T(n,k)-(n-k)/n=k(2k-n-1)/(n(n-1))<0\). So the population is two-valued, the
largest atom of \(X\) is a largest atom of \(\Hyp(n,i,k)\), and
Lemma~\ref{lem:atom} classifies equality. At half-sampling the three-value
bound is \(T(2k,k)=(k-1)/(2k-1)<q_k\). For two-valued populations the largest
atom of \(\Hyp(2k,i,k)\) is its central, modal atom: for odd \(i\) every atom
is at most \(1/2\) by Theorem~\ref{thm:parity}, and for even \(i=2s\) the
central atom is \(c_s\), which equals \(q_k\) only for \(s\in\{1,k-1\}\)
(Section~\ref{sec:small}).
\end{proof}

The case \(n=2k\) of this bound is due to \cite{MVV}. In particular the
minimum min-entropy \(H_\infty=-\log_2\max_tp_t\) is minus the binary logarithm
of the displayed bound; at \(n=2r+1\ge5\) it has more minimisers than any
finite order.

The chain in Lemma~\ref{lem:chain} moves by at least the gap \(g\) at each
step, so the atom bound extends to short intervals. This is a small-ball
(Littlewood--Offord type) inequality for sampling without replacement, at the
natural scale of the population's largest gap.

\begin{corollary}[Small-ball inequality]\label{cor:smallball}
Let \(x\) be a nonconstant real population and let \(g\) be the largest gap
between consecutive distinct values of \(x\). Then every half-open interval
\(I\) of length at most \(g\) satisfies \(\Prb(X\in I)\le\beta(n,k)\), where
\(\beta(n,k)\) is the right side of Corollary~\ref{cor:concentration}. More
generally, a union of \(\ell\) such intervals has probability at most the sum
of the \(\ell\) largest entries of the corresponding frontier vector of
Theorem~\ref{thm:frontier}; at \(n=2k\ge4\) this means the larger of the
sums of the \(\ell\) largest entries of \(U\) and of \(Q_k\). The bounds are
attained by two-valued populations.
\end{corollary}
\begin{proof}
Apply Lemma~\ref{lem:chain} at the largest gap (it holds for every
\(k\); the bound for \(\Hyp(n,i,k)\) with \(k>n/2\) follows from the case
\(n-k\) because the two laws have the same masses). Consecutive chain sums differ
by at least \(g\), so a half-open interval of length at most \(g\) contains at
most one of them, and a union of \(\ell\) such intervals at most \(\ell\). The
argument in the proof of Lemma~\ref{lem:chain} then bounds the probability by
the sum of the \(\ell\) largest masses of \(\Hyp(n,i,k)\), and
Lemma~\ref{lem:atom} and Theorem~\ref{thm:parity} bound these by the frontier.
For a two-valued population the bound with \(\ell=1\) is the largest atom,
which Corollary~\ref{cor:concentration} attains.
\end{proof}

For a family \(\mathcal F\subseteq\Omega_{n,k}\), let \(M(\mathcal F)\) be its
real incidence matrix: rows are the sets in \(\mathcal F\), columns the
elements of \([n]\).

\begin{corollary}[One unit of rank deficiency]\label{cor:rank}
The largest families with \(\rank M(\mathcal F)<n\) have size
\[
 \binom{n-1}{k}\ (k<n/2),\qquad
 2\binom{n-2}{k-1}\ (n=2k\ge4),\qquad
 \binom{n-1}{k-1}\ (k>n/2),\qquad
 1\ (n=2,\ k=1).
\]
\end{corollary}
\begin{proof}
A nonzero null vector \(x\) of \(M(\mathcal F)\) makes \(L_x\) vanish on
\(\mathcal F\). If \(x\) is constant the family is empty; otherwise
Corollary~\ref{cor:concentration} bounds \(|\mathcal F|\). For \(k<n/2\) all
\(k\)-sets avoiding one element attain the bound; for \(k>n/2\) all \(k\)-sets
containing one element do, with null vector \((-(k-1),1,\ldots,1)\). At
half-sampling take all \(k\)-sets containing exactly one of two fixed elements,
with null vector \(((k-1)^2,(-1)^{2k-2})\). The case \(n=2\) is immediate.
\end{proof}

Corollary~\ref{cor:rank} is not new: it is the case \(s=t=1\) of the
rank-extremal problem studied by Grosu, Person and Szab\'o \cite{GPS}, who
report that this case was determined by Odlyzko and that Ahlswede, Aydinian
and Khachatrian solved it for ordinary incidence matrices and every rank
deficiency
\cite[Section~2, Remark]{GPS}. It is included because it follows in a few
lines from Corollary~\ref{cor:concentration}.

\section{The hypergeometric tail problem}\label{sec:tail}

This section proves Theorem~\ref{thm:tail} and states its consequence for the
MMS problem. The proof uses one external lemma.

\begin{lemma}[{Ai, Karagiannis, Ku\v{z}elka and Pelekis \cite[Theorem~2]{AKKP}}]\label{lem:AKKP}
Let \(Y\) be an integer-valued random variable with \(\E|Y|<\infty\) whose
survival probabilities \(\Prb(Y\ge j)\) form a log-concave sequence. If
\(m\) is an integer with \(\E Y\ge m\) and
\(\Prb(Y=m+1)\ge\tfrac12\Prb(Y=m)>0\), then \(\Prb(Y\ge m+1)\ge1/4\).
\end{lemma}
(The cited theorem allows any ratio \(\gamma\in(0,1/2]\) in place of
\(\tfrac12\), with bound \(\gamma/(1+2\gamma)\).) A hypergeometric law has
log-concave masses, since
\[
 \frac{\Prb(H=j+1)}{\Prb(H=j)}=\frac{(i-j)(k-j)}{(j+1)(n-i-k+j+1)}
\]
decreases in \(j\) on its support, which is an interval; log-concave masses
give log-concave survival probabilities, so Lemma~\ref{lem:AKKP} applies to
\(H\sim\Hyp(n,i,k)\).

\begin{proof}[Proof of Theorem~\ref{thm:tail}]
Let \(n\ge4k\), \(H\sim\Hyp(n,i,k)\) and \(\mu=ik/n\). Since \(H\) is
integer-valued, \(\Prb(H\ge\mu)=\Prb(H\ge\lceil\mu\rceil)\).

\emph{Case A: \(\mu\le1\).} Then \(\lceil\mu\rceil=1\) and
\(\Prb(H\ge1)=1-\binom{n-i}{k}/\binom nk\). For \(i=1\) this equals \(k/n\).
For \(i\ge2\), \(\binom{n-i}{k}\le\binom{n-2}{k}<\binom{n-1}{k}\), since
\(\binom{n-1}{k}-\binom{n-2}{k}=\binom{n-2}{k-1}>0\); so \(\Prb(H\ge1)>k/n\).

\emph{Case B: \(\mu>1\) and \(\mu\notin\Z\).} Here
\(\lambda=(n-i)(n-k)/n\ge(n-k)/n\ge3/4\), so both \(\mu\) and \(\lambda\)
exceed \(\log(4/3)\), and \cite[Theorem~3(1)]{AKKP} gives
\(\Prb(H\ge\mu)=\Prb(H>\mu)\ge1/4\ge k/n\). For completeness we recall their
argument in the present setting. If \(i=n-1\), then
\(\lceil\mu\rceil=k\) and \(\Prb(H=k)=(n-k)/n\ge3/4\). Let \(i\le n-2\) and
\(m=\lfloor\mu\rfloor\ge1\). Because \(\mu\) lies strictly between the integers
\(m\) and \(m+1\) and the support of \(H\) is an integer interval containing
\(\mu\) in its convex hull, both \(m\) and \(m+1\) carry positive mass. Put
\(\lambda=n-i-k+\mu=(n-i)(n-k)/n\ge2(n-k)/n\ge3/2\); then the integer
\(n-i-k+m=\lambda-(\mu-m)\) exceeds \(1/2\), so \(n-i-k+m\ge1\). Also
\(ik-nm\ge1\). The identity of \cite[Section~3]{AKKP}
\[
 2(i-m)(k-m)-(m+1)(n-i-k+m+1)=(m-1)(n-i-k+m-1)+2(ik-nm-1)
\]
shows that the left side is nonnegative, i.e.
\(\Prb(H=m+1)\ge\tfrac12\Prb(H=m)\). Lemma~\ref{lem:AKKP} gives
\(\Prb(H\ge\mu)=\Prb(H\ge m+1)\ge1/4\ge k/n\).

\emph{Case C: \(\mu=j\in\Z\), \(j\ge2\).} If \(j-1\) is outside the support,
then \(\Prb(H\ge j)=1\). Otherwise, with the cell counts
\(a=j\), \(b=i-j\), \(c=k-j\), \(d=n-i-k+j\) of Section~\ref{sec:support},
all positive, \eqref{eq:ratios} gives
\(\Prb(H=j)/\Prb(H=j-1)=(b+1)(c+1)/(bc)>1\). Lemma~\ref{lem:AKKP} with
\(m=j-1\) gives \(\Prb(H\ge j)\ge1/4\ge k/n\).

So \(\Prb(H\ge\mu)\ge k/n\) for all \(i\), with equality at \(i=1\). In Cases
B and C the bound is at least \(1/4\), which exceeds \(k/n\) when \(n>4k\), and
in Case A it is strict for \(i\ge2\). Hence for \(n>4k\) only \(i=1\) attains
\(k/n\).
\end{proof}

The proof shows slightly more: for \(n\ge4k\) every \(i\) with
\(ik/n>1\) has tail at least \(1/4\). The finite check
described at the end of the paper finds, for all \(k\le40\) and \(4k\le n\le10k\), that
every \(i\ge2\) has tail at least \(\tfrac54\cdot\tfrac kn\); so in that range
\(i=1\) is also the unique minimiser at \(n=4k\).

\begin{corollary}[Aydinian--Blinovsky implies MMS]\label{cor:AB}
For \(n\ge4k\), \(\min_{1\le i\le n-1}\Prb(\Hyp(n,i,k)\ge ik/n)=k/n\).
Consequently, if the conjecture of Aydinian and Blinovsky \cite{AB} holds,
then every real population with zero total and \(n\ge4k\) satisfies
\(\Prb(X\ge0)\ge k/n\), which is the MMS conjecture.
\end{corollary}
\begin{proof}
The first statement is Theorem~\ref{thm:tail}. In the conjecture of \cite{AB},
as formulated in \cite{AKP,AP}, the extremal populations are \(i\) copies of
\(1/i\) and \(n-i\) copies of \(-1/(n-i)\); for them \(X\ge0\) exactly when
\(\Hyp(n,i,k)\ge ik/n\).
\end{proof}

\begin{remark}[Below \(n=4k\)]\label{rem:below4k}
The restriction \(n\ge4k\) cannot simply be dropped. Ai, Ku\v{z}elka and
Pelekis note that \(i=1\) is not minimal at \(n=3k+1\) for \(k\ge3\) \cite{AKP}. The finite
check described at the end of the paper shows that for each \(2\le k\le40\) there are
values \(2k<n<4k\) at which some \(i\ge2\) has a smaller tail than \(i=1\), the
largest such ratio \(n/k\) being \(10/3\) (at \(k=3\), \(n=10\)); for these \(k\)
there is none with \(10/3<n/k<4\). The exact range of \(n/k\) where the single outlier is optimal
remains open.
\end{remark}

\section{Abelian groups}\label{sec:abelian}

This section proves Theorem~\ref{thm:groups}. The real proofs used order; over
a general Abelian group only cancellation is available, and parity takes the
place of order. Throughout, \(f(S)=c+\sum_{i\in S}a_i\) with values in \(G\),
and \(\delta_G(n,k)\) is the smallest proportion of \(k\)-sets with
\(f(S)\ne0\), over \(f\) not identically zero on \(\Omega_{n,k}\). (A formal
expression with nonzero coefficients may still vanish identically on the
slice; such \(f\) are excluded.)

\begin{lemma}[Independent sets in unimodal sequences]\label{lem:unimodal-parity}
Let \(w_0,\ldots,w_m\) be nonnegative and unimodal. Then the largest total
weight of a set \(I\subseteq\{0,\ldots,m\}\) containing no two consecutive
integers is \(\max\{\sum_{j\ \mathrm{even}}w_j,\sum_{j\ \mathrm{odd}}w_j\}\).
\end{lemma}
\begin{proof}
Both parity classes are admissible, so the maximum is at least the right side.
Conversely, let \(I\) be admissible and let \(h\) be a mode. Call consecutive
selected indices \(a<b\) a \emph{parity change} if \(b-a\) is odd; then
\(b-a\ge3\). If \(a<h\), shift every selected index \(\le a\) one step to the
right: these indices lie on the weakly increasing arm \(\{0,\ldots,h\}\), so
no weight decreases, and the gap \(b-a\) drops to an even number \(\ge2\). If
\(a\ge h\), then every selected index \(\ge b\) exceeds \(h\) and lies on the
weakly decreasing arm; shift each of them one step to the left, with the same
effect. In both cases all other gaps keep their length, the set stays
admissible, and the number of parity changes drops by one. After finitely many
steps \(I\) lies in one parity class, without loss of weight, and
nonnegativity bounds its weight by that class.
\end{proof}

For the binary Krawtchouk polynomials
\(K_k^{(n)}(i)=\sum_j(-1)^j\binom ij\binom{n-i}{k-j}\) we use the following
classical bound.

\begin{lemma}[{Dumer--Kapralova \cite{DK}; see \cite[Theorem~1.2]{BCIM}}]\label{lem:known-krawtchouk}
For \(1\le i\le n-1\),
\(|K_k^{(n)}(i)|\le|K_k^{(n)}(1)|\) if \(k\ne n/2\), and
\(|K_k^{(n)}(i)|\le|K_k^{(n)}(2)|\) if \(k=n/2\). Moreover
\[
 \frac{K_k^{(n)}(1)}{\binom nk}=\frac{n-2k}{n},\qquad
 \frac{K_k^{(n)}(2)}{\binom nk}=\frac{(n-2k)^2-n}{n(n-1)} .
\]
\end{lemma}
Brouwer, Cioab\u{a}, Ihringer and McGinnis state the inequalities as their
Theorem~1.2, attributing them to \cite[Corollary~10]{DK}, and give a
self-contained proof; the evaluations follow from the definition. In the
balanced case we use the lemma only for \(n\ge4\).

\begin{proof}[Proof of Theorem~\ref{thm:groups}: distance]
By complementation assume \(k\le n/2\). If all \(a_i\) are equal, \(f\) is
constant on the slice and, being not identically zero, has full support. If
three of the \(a_i\) are distinct, Lemma~\ref{lem:three} gives a proportion at
least \(2k(n-k)/(n(n-1))\ge k/n\), which suffices in both regimes (at \(n=2k\),
\(2k^2/(2k(2k-1))=k/(2k-1)>(k-1)/(2k-1)\)).

It remains to treat two values, with multiplicities \(i\) and \(n-i\). Then
\(f(S)=A+gJ\) with \(A\in G\), \(g\ne0\) and \(J\sim\Hyp(n,i,k)\). The set
of integers \(j\) with \(A+gj=0\) contains no two consecutive integers, since
otherwise \(g=0\). The masses \(w_j=\Prb(J=j)\) are unimodal (they are
log-concave, Section~\ref{sec:tail}), so Lemma~\ref{lem:unimodal-parity} gives
\[
 \Prb(f=0)\le\max\{\Prb(J\ \mathrm{even}),\Prb(J\ \mathrm{odd})\}
 =\frac12\Bigl(1+\frac{|K_k^{(n)}(i)|}{\binom nk}\Bigr),
\]
and Lemma~\ref{lem:known-krawtchouk} bounds this by \((n-k)/n\) when
\(n>2k\) and by \(k/(2k-1)\) when \(n=2k\ge4\). For \((n,k)=(2,1)\) a function
on the two-point slice that is not identically zero has support at least one.

For attainment, let \(g\ne0\). The function \(g\one_{\{1\in S\}}\) has
nonzero proportion \(k/n\) (and \(1/2\) at \(n=2\)). At balance,
\(f(S)=g(\one_{\{1\in S\}}+\one_{\{2\in S\}}-1)\) vanishes exactly when one of
\(1,2\) lies in \(S\), and otherwise equals \(\pm g\ne0\), whatever the order of
\(g\); its nonzero proportion is \((k-1)/(2k-1)\).
\end{proof}

Consequently, if \(a_1,\ldots,a_n\in G\) are not all equal, every atom of
\(X=\sum_{i\in S}a_i\) has mass at most \(1-\delta_G(n,k)\): two \(k\)-sets that
share \(k-1\) elements and differ at unequal coefficients have different sums,
so \(X-z\) is not identically zero for any \(z\in G\).

\begin{proof}[Proof of Theorem~\ref{thm:groups}: entropy]
Write \(G[2]=\{g\in G:2g=0\}\). For \(n\ne2k\), every atom has mass at most
\(1-r/n\), so \(\Ber(k/n)\), sorted as \((1-r/n,r/n)\), majorizes the law;
\((g,0,\ldots,0)\) with \(g\ne0\) attains it. For \(n=2\) the law is \(U\).

Let \(n=2k\ge4\). Every atom has mass at most \(q_k\), so \(\Ber(q_k)\)
majorizes every law. If some \(g\ne0\) has \(2g=0\), the population
\((g,g,0,\ldots,0)\) has sum \(g\) when exactly one of the first two entries
is chosen (probability \(q_k\)) and \(0\) otherwise; this is \(\Ber(q_k)\).

If \(G[2]=\{0\}\), put \(T=\sum_ia_i\). Complementation gives
\(X\overset d=T-X\), and the involution \(y\mapsto T-y\) has at most one fixed
point, since two fixed points would differ by a nonzero element of \(G[2]\).
If every atom has mass at most \(1/2\), the law is majorized by \(U\).
Otherwise an atom of mass \(c>1/2\) must be the fixed point, with \(c\le q_k\);
every other atom has a distinct partner of equal mass, so the second largest
mass is at most \((1-c)/2\) and the two largest sum to at most
\((1+q_k)/2\). Hence the law is majorized by \(Q_k\). The populations
\((g,0,\ldots,0)\) and \((g,g,0,\ldots,0)\) attain \(U\) and \(Q_k\), since
\(0,g,2g\) are distinct. Schur-concavity and Section~\ref{sec:entropy} finish
the proof.
\end{proof}

\begin{remark}[Scope]\label{rem:Zq}
The case \(G[2]=\{0\}\) includes torsion-free groups and finite groups of odd
order. Theorem~\ref{thm:groups} concerns the affine problem. The zero-sum
problem of Theorem~\ref{thm:support} does \emph{not} carry over to finite
groups: over \(\Z_2\), the vector \((0^7,1^4)\) has zero total and only \(224\)
nonzero \(5\)-sums among \(\binom{11}5\), fewer than the real minimum
\(2\binom94=252\); over \(\Z_3\), \((1^9,2^3)\) with \(k=5\) gives \(414<420\);
over \(\Z_5\), \((0,1^{10})\) with \(k=5\) gives \(210<252\). The reason is
that the zero event of a two-valued population becomes a residue class of the
hypergeometric count rather than a single atom, which defeats
Lemma~\ref{lem:meanatom}. Determining the zero-sum minimum over \(\Z_q\) is
open. No classification of the group-valued entropy minimisers is attempted.
\end{remark}

\section{Open problems}\label{sec:open}

\begin{enumerate}
\item \emph{The sign problem below \(4k\), and the reduction itself.} Determine
the exact range of \(n/k\) in which the single outlier minimises
\(\Prb(\Hyp(n,i,k)\ge ik/n)\) (Remark~\ref{rem:below4k}); and prove the
Aydinian--Blinovsky reduction, which with Theorem~\ref{thm:tail} would give the
MMS conjecture.
\item \emph{Higher Johnson eigenspaces.} Find the minimum support of the
eigenspace of the \(\ell\)-th eigenvalue of \(J(n,k)\) for all \(n\), beyond
the large-\(n\) result of \cite{VMV}; or the analogous first-eigenspace problem
for Grassmann graphs.
\item \emph{Finite groups.} Determine the zero-sum minimum support over
\(\Z_q\) (Remark~\ref{rem:Zq}) and classify the group-valued entropy
minimisers.
\item \emph{Higher degree.} Find the exact distance of degree-\(d\) polynomials
on the Boolean slice for \(d\ge2\), where \cite{ABSS} gives approximate bounds.
\end{enumerate}

\section*{Data, verification and AI use}\label{sec:data}
\addcontentsline{toc}{section}{Data, verification and AI use}
All proofs above are complete as written and use no computer step. The
accompanying scripts check the statements exactly (integer or rational
arithmetic) on finite ranges: an exhaustive search over all candidate vectors
for \(n\le7\), which is complete because every minimiser spans a line cut out
by \(0/1\) equations and so has integer entries bounded by the maximal
determinant of a \(0/1\) matrix of order \(n-1\); the support formula on
integer grids and the integer-mean comparison through \(n=300\); the atom and
parity majorization statements; all populations with entries in \(\{0,1,2,3\}\) and \(n\le10\) against the entropy
and equality statements; the hypergeometric tail statements for all
\(k\le40\), \(2k<n\le10k\); the small-ball inequality on random rational
populations; and the group statements on small Abelian groups. The scripts and their outputs are part of the
release. This paper was prepared with substantial assistance from AI
systems, which produced the proofs, computations and text; it has been checked
by AI-assisted internal review and exact computation, not by external referees
or a proof assistant.

\begin{thebibliography}{99}
\bibitem{AB}
H.~Aydinian and V.~Blinovsky,
A remark on the problem of nonnegative \(k\)-subset sums,
\emph{Problems of Information Transmission} \textbf{48}(4) (2012), 347--351.
\url{https://doi.org/10.1134/S0032946012040059}.

\bibitem{ABSS}
P.~Amireddy, A.~R.~Behera, S.~Srinivasan and M.~Sudan,
A near-optimal polynomial distance lemma over Boolean slices,
in \emph{52nd International Colloquium on Automata, Languages, and
Programming (ICALP 2025)}, LIPIcs \textbf{334} (2025), 11:1--11:17.
\url{https://doi.org/10.4230/LIPIcs.ICALP.2025.11}.
Full version (Theorem~6.1): ECCC TR25-056,
\url{https://eccc.weizmann.ac.il/report/2025/056/}.

\bibitem{AKP}
J.~Ai, O.~Ku\v{z}elka and C.~Pelekis,
An entropic analogue of the MMS conjecture,
preprint, arXiv:2606.30486v1 (2026).
\url{https://arxiv.org/abs/2606.30486v1}.

\bibitem{AKP25}
J.~Ai, O.~Ku\v{z}elka and C.~Pelekis,
Concentration inequalities for the sum in sampling without replacement: an
approach via majorization,
preprint, arXiv:2503.20473 (2025).
\url{https://arxiv.org/abs/2503.20473}.

\bibitem{AKKP}
J.~Ai, V.~Karagiannis, O.~Ku\v{z}elka and C.~Pelekis,
On the probability that a discrete IFR random variable exceeds its mean,
preprint, arXiv:2610.24726v1 (2026).
\url{https://arxiv.org/abs/2610.24726v1}.

\bibitem{AP}
J.~Ai and C.~Pelekis,
On lower bounds for hypergeometric tails,
\emph{Statistics \& Probability Letters} \textbf{236} (2026), 110798.
\url{https://doi.org/10.1016/j.spl.2026.110798}.
Preprint: arXiv:2601.09485.

\bibitem{BCIM}
A.~E.~Brouwer, S.~M.~Cioab\u{a}, F.~Ihringer and M.~McGinnis,
The smallest eigenvalues of Hamming graphs, Johnson graphs and
other distance-regular graphs with classical parameters,
\emph{Journal of Combinatorial Theory, Series B} \textbf{133} (2018), 88--121.
\url{https://doi.org/10.1016/j.jctb.2018.04.005}.

\bibitem{DK}
I.~Dumer and O.~Kapralova,
Spherically punctured biorthogonal codes,
\emph{IEEE Transactions on Information Theory} \textbf{59}(9) (2013),
6010--6017.
\url{https://doi.org/10.1109/TIT.2013.2250579}.

\bibitem{GPS}
C.~Grosu, Y.~Person and T.~Szab\'o,
On the rank of higher inclusion matrices,
\emph{Journal of the London Mathematical Society} \textbf{90}(2) (2014), 333--349.
\url{https://doi.org/10.1112/jlms/jdu029}.

\bibitem{Levin}
B.~Levin,
A stochastic ordering of multiple hypergeometric laws: peakedness of category
counts about half the population category sizes is symmetric unimodal in the
sample size,
\emph{Sequential Analysis} \textbf{43}(4) (2024), 417--431.
\url{https://doi.org/10.1080/07474946.2024.2379901}.
Preprint: arXiv:2311.17884.

\bibitem{MS88}
N.~Manickam and N.~M.~Singhi,
First distribution invariants and EKR theorems,
\emph{Journal of Combinatorial Theory, Series A} \textbf{48}(1) (1988), 91--103.
\url{https://doi.org/10.1016/0097-3165(88)90077-5}.

\bibitem{MSch}
L.~Mattner and J.~Schulz,
On normal approximations to symmetric hypergeometric laws,
\emph{Transactions of the American Mathematical Society} \textbf{370}(1) (2018), 727--748.
\url{https://doi.org/10.1090/tran/6986}.

\bibitem{MVV}
I.~Yu.~Mogilnykh, K.~V.~Vorob'ev and A.~A.~Valyuzhenich,
MMS-type problems for Johnson scheme,
\emph{Siberian Electronic Mathematical Reports} \textbf{15} (2018), 1663--1670.
\url{https://doi.org/10.33048/semi.2018.15.137}.
Preprint: arXiv:1810.04075.

\bibitem{MWW}
M.~Madiman, L.~Wang and J.~O.~Woo,
Majorization and R\'enyi entropy inequalities via Sperner theory,
\emph{Discrete Mathematics} \textbf{342}(10) (2019), 2911--2923.
\url{https://doi.org/10.1016/j.disc.2019.03.002}.

\bibitem{Pok}
A.~Pokrovskiy,
A linear bound on the Manickam--Mikl\'os--Singhi conjecture,
\emph{Journal of Combinatorial Theory, Series A} \textbf{133} (2015), 280--306.
\url{https://doi.org/10.1016/j.jcta.2015.01.007}.

\bibitem{SV}
E.~Sotnikova and A.~Valyuzhenich,
Minimum supports of eigenfunctions of graphs: a survey,
\emph{The Art of Discrete and Applied Mathematics} (2021).
\url{https://doi.org/10.26493/2590-9770.1404.61e}.
Preprint: arXiv:2102.11142.

\bibitem{VMV}
K.~Vorob'ev, I.~Mogilnykh and A.~Valyuzhenich,
Minimum supports of eigenfunctions of Johnson graphs,
\emph{Discrete Mathematics} \textbf{341}(8) (2018), 2151--2158.
\url{https://doi.org/10.1016/j.disc.2018.04.018}.
\end{thebibliography}
\end{document}
